\section*{अध्याय \ref{ch:b3:canonical-ensemble} --- विहित समुदाय}

\begin{solution}{exo:b3:canonical-ensemble:1}
(क) $\mathcal P_2/\mathcal P_1 = (g_2/g_1)\eu^{-\beta\epsilon}$।
(ख) $\beta\epsilon = 2.1/(8.62\times10^{-5} \times 2500) = 9.7$:
अंश $3\eu^{-9.7} \approx \num{2e-4}$। (ग) कोई लौ
$\sim\num{e15}$ सोडियम परमाणु प्रति cm$^3$ ढोती है: उनमें से $10^{-4}$ भी,
हर कुछ नैनोसेकंड में चक्र लगाते हुए, प्रति सेकंड $10^{19}$ पीले
फ़ोटॉन उँडेल देते हैं --- आँखें चौंधियाने भर। (घ) $3\eu^{-\beta\epsilon} = 0.1$:
$T = \epsilon/(k_{\text{B}}\ln30) \approx \qty{7200}{K}$ --- यह
किसी तारे का प्रकाशमंडल है, कोई लौ नहीं।
\end{solution}

\begin{solution}{exo:b3:canonical-ensemble:2}
(क) $z = 1 + \eu^{-\beta\epsilon} + \eu^{-2\beta\epsilon}$। (ख)
$\langle E\rangle = \epsilon(\eu^{-\beta\epsilon} +
2\eu^{-2\beta\epsilon})/z$: $\to 0$ निम्न $T$ पर, और उच्च पर
$\to \epsilon$ (माध्य स्तर)। (ग) $\mathcal P_1 = 1/(\eu^{\beta
\epsilon} + 1 + \eu^{-\beta\epsilon})$ $T$ के साथ एकदिष्ट रूप से
बढ़ता है और अपने उच्चतम $1/3$ तक केवल $T \to \infty$ पर पहुँचता है। (घ)
बोल्ट्ज़मान भार $T > 0$ पर ऊर्जा के साथ केवल घटते हैं; और
जनसंख्या-प्रतिलोमन को
\cref{ex:b3:microcanonical-ensemble:negative} वाले \omterm{ex:b3:microcanonical-ensemble:negative}{ऋणात्मक ताप}
चाहिए, जिन तक कोई भंडार पहुँच ही नहीं सकता।
\end{solution}

\begin{solution}{exo:b3:canonical-ensemble:3}
(क) $E_p = mgh$ पर बोल्ट्ज़मान गुणक, एकसमान $T$ पर। (ख) $h_0 =
k_{\text{B}}T/mg = \qty{8.4}{km}$। (ग) $\eu^{-8848/8440} \approx
0.35$: एक तिहाई वायुमंडल --- इसीलिए शिखर-यात्री ऑक्सीजन ढोते हैं।
(घ) असली वायुमंडल समतापी नहीं है: ताप ऊँचाई के साथ गिरता है, अतः
प्रोफ़ाइल किसी एक चरघातांकी से हटकर मुड़ जाता है।
\end{solution}

\begin{solution}{exo:b3:canonical-ensemble:4}
(क) तीन स्थानांतरण: $\tfrac32 R = \qty{12.5}{J/(mol.K)}$,
ठीक-ठीक सहमत। (ख) दो घूर्णन जोड़िए: $\tfrac52 R =
\qty{20.8}{J/(mol.K)}$, जैसा नापा गया: कंपन जमा हुआ है। (ग)
$\tfrac72 R = \qty{29.1}{J/(mol.K)}$ की भविष्यवाणी; और मापा गया
$\approx26$ दिखाता है कि कंपन केवल आंशिक रूप से पिघला है ($\theta_{\text{कं}}
\approx \qty{3400}{K}$)। (घ) $3R = \qty{24.9}{J/(mol.K)}$:
अर्थात् द्यूलों--प्टी, जिसे कमरे के ताप पर अधिकतर धातुएँ मानती हैं
और हीरा तोड़ देता है --- वही जमाव की कहानी
(\cref{exo:b3:canonical-ensemble:6})।
\end{solution}

\begin{solution}{exo:b3:canonical-ensemble:5}
(क) $z = \eu^{-\beta\hbar\omega/2}/(1 - \eu^{-\beta\hbar\omega})$।
(ख) $\langle E\rangle = -\partial_\beta\ln z$ कथित रूप देता है,
जहाँ $\langle n\rangle = 1/(\eu^{\beta\hbar\omega} - 1)$। (ग)
$C \to k_{\text{B}}$ उच्च $T$ पर; और निम्न पर $C \approx k_{\text{B}}(\beta
\hbar\omega)^2\eu^{-\beta\hbar\omega} \to 0$। (घ) जहाँ
$\hbar\omega$ प्रकाश के किसी क्वांटम की ऊर्जा हो, वहाँ $\langle n\rangle$
प्रति मोड तापीय फ़ोटॉन-संख्या है: प्लांक का नियम एक ही अध्याय दूर है।
\end{solution}

\begin{solution}{exo:b3:canonical-ensemble:6}
(क) संख्यात्मक रूप से, $C = \tfrac12 \times 3Nk_{\text{B}}$
$T/\theta_{\text{E}} \approx 0.34$ के पास। (ख) $\theta_{\text{E}}/T =
4.3$: $C/3Nk_{\text{B}} = 4.3^2\eu^{-4.3}/(1 - \eu^{-4.3})^2
\approx 0.26$ --- कमरे के ताप पर हीरे के पास चिरसम्मत ऊष्माधारिता
का मुश्किल से एक चौथाई है: वह ``विषमता'' खुले में खड़ा क्वांटम
जमाव है। (ग) $\theta_{\text{E}} = \qty{90}{K}$
सीसे को \qty{300}{K} पर गहरे चिरसम्मत प्रांत में रखता है। (घ)
एक-आवृत्ति वाली मान्यता: असली ठोसों में मोडों का पूरा वर्णक्रम होता
है, जो लंबी-तरंगदैर्घ्य वाली ध्वनि तक जाता है, और जिनके सस्ते
क्वांटा $T^3$ नियम देते हैं --- डिबाई की मरम्मत,
\cref{ch:b3:photons-phonons} में।
\end{solution}

\begin{solution}{exo:b3:canonical-ensemble:7}
(क) यहाँ $f(v) \propto v^2\eu^{-mv^2/2k_{\text{B}}T}$: $v_{\text{p}} =
\sqrt{2k_{\text{B}}T/m}$: N$_2$ के लिए \qty{413}{m/s},
और H$_2$ के लिए \qty{1540}{m/s}। (ख) $mv_{\text{पल}}^2/2k_{\text{B}}T
\approx 730$ N$_2$ के लिए, और $52$ H$_2$ के लिए। (ग) $\eu^{-730}$
कभी नहीं है; और $\eu^{-52} \sim \num{e-23}$ प्रति निवास-काल धीमा है
--- पर तप्त ऊपरी वायुमंडल ($\sim\qty{1000}{K}$) हाइड्रोजन का
घातांक नरम करके $\sim15$ कर देता है: भूवैज्ञानिक काल में हाइड्रोजन
रिस जाता है, और नाइट्रोजन टिका रहता है। (घ) पलायन वेग और बहिर्मंडल
का ताप: बृहस्पति का \qty{60}{km/s} कूप हाइड्रोजन तक का
घातांक खगोलीय बना देता है --- गैस-दानव वही रखते हैं जो छोटे गरम लोक
खो देते हैं।
\end{solution}

\begin{solution}{exo:b3:canonical-ensemble:8}
(क) $\partial_\beta^2\ln Z = \langle E^2\rangle - \langle
E\rangle^2$, और $C = \partial_T\langle E\rangle =
-k_{\text{B}}\beta^2\partial_\beta\langle E\rangle$। (ख) दोनों
पक्ष $Nk_{\text{B}}(\beta\epsilon)^2\eu^{\beta\epsilon}/(
\eu^{\beta\epsilon} + 1)^2$ देते हैं। (ग) $\langle E\rangle \propto N$,
और $\Delta E \propto \sqrt N$। (घ) कोई तंत्र गरम किए जाने पर कितनी
प्रबलता से अनुक्रिया करता है, यह उतना ही है जितना उसकी ऊर्जा
स्वतःस्फूर्त रूप से डगमगाती है --- अनुक्रिया और उतार-चढ़ाव उसी एक
द्वितीय अवकलज के दो पाठ हैं, और यह ढाँचा
(उतार-चढ़ाव--अपव्यय) पूरी भौतिकी में लौटता रहता है।
\end{solution}

\begin{solution}{exo:b3:canonical-ensemble:9}
(क) $\ln2 = E_a\,\Delta T/k_{\text{B}}T^2$, जहाँ $\Delta T =
\qty{10}{K}$, $T = \qty{298}{K}$: $E_a =
0.693\,k_{\text{B}}T^2/10 \approx \qty{0.55}{eV}$। (ख)
$298$ से \qty{278}{K} तक: गुणक $\eu^{E_a(1/278 - 1/298)/k_{\text{B}}}
\approx 4.4$ जितना धीमा --- इसीलिए शीतक भोजन बचाते हैं। (ग)
$373 \to \qty{366}{K}$: दर $\approx1.4$ से गिरती है: ग्यारह मिनट
वाले अंडे को पौन घंटा चाहिए। (घ) अभिक्रियाएँ रोधिका से ऊपर वाले
अणुओं की चरघातांकी पूँछ जीतती है: ताप ज़रा-सा खिसकाइए और उस पूँछ की
जनसंख्या ज़बरदस्त खिसक जाती है --- माध्य से शायद ही कुछ फ़र्क पड़ता है।
\end{solution}

\begin{solution}{exo:b3:canonical-ensemble:10}
(क) $f_{\text{मु}} = 1/(1 + \eu^{-\beta(\Delta E -
T\Delta S)})$, जहाँ खुली अवस्था की एन्ट्रॉपी उसकी
\omterm{prop:b3:canonical-ensemble:machine}{मुक्त ऊर्जा} में समा गई है। (ख) $T_{\text{m}} = \Delta E/\Delta S$ पर
दोनों मुक्त ऊर्जाएँ बराबर हो जाती हैं: आधा-आधा। (ग) $T_{\text{m}} =
\num{4.8e-19}/\num{1.38e-21} = \qty{348}{K}$ (\qty{75}{\celsius});
चौड़ाई $\delta T \sim k_{\text{B}}T_{\text{m}}^2/\Delta E \approx
\qty{3.5}{K}$। (घ) सहकारिता एक साथ टूटते संपर्कों की संख्या से
$\Delta E$ और
$\Delta S$ दोनों को गुणा कर देती है, जिससे
$T_{\text{m}}$ वही रहता है पर चौड़ाई $\propto 1/\Delta E$ सिकुड़ जाती
है: DNA की लड़ियाँ दो-एक केल्विन में अलग हो जाती हैं, और यही किसी
PCR मशीन को ``पिघला'' और ``जुड़ा'' के बीच स्वच्छ चक्र लगाने देता है।
\end{solution}

\begin{solution}{exo:b3:canonical-ensemble:11}
(क) $z = 2\cosh(\beta\mu B)$; $M = N\mu\tanh(\beta\mu B) \approx
N\mu^2B/k_{\text{B}}T$ छोटे तर्क के लिए: अर्थात् क्यूरी का $1/T$। (ख)
$\mu_{\text{B}}B/k_{\text{B}}T$: \qty{300}{K} पर $\num{2.2e-3}$;
और \qty{1}{K} पर $0.67$ --- उदासीन से लेकर प्रबल रूप से संरेखित तक। (ग)
$S$ केवल $\mu B/k_{\text{B}}T$ पर निर्भर करता है; अतः स्थिर
एन्ट्रॉपी पर $B$ घटाना $T$ को उसी अनुपात में नीचे खींचता है: $\qty{1}{K} \to
\qty{0.1}{K}$। (घ) जब $k_{\text{B}}T$ चक्रण--चक्रण ऊर्जा तक पहुँच
जाए, तब एन्ट्रॉपी क्षेत्र के वश में नहीं रहती:
$\mu_0\mu_{\text{B}}^2/4\pi a^3 \approx \qty{7e-26}{J} \sim
\qty{5}{mK}$ --- अर्थात् चिरपरिचित तली (नाभिकीय आघूर्ण, जो हज़ार
गुना दुर्बल हैं, उसे माइक्रोकेल्विन तक धकेल देते हैं)।
\end{solution}

\begin{solution}{exo:b3:canonical-ensemble:12}
(क) $z_{\text{घू}} = \sum_J(2J+1)\eu^{-\beta BJ(J+1)}
\to \int(2J+1)\eu^{-\beta BJ(J+1)}\dd J = k_{\text{B}}T/B$: तब
$\langle E\rangle = k_{\text{B}}T$ और $C_{\text{घू}} =
k_{\text{B}}$। (ख) H$_2$: $\theta_{\text{घू}} = \qty{87}{K}$ ---
पायदान प्रयोगशाला के परास में बैठता है; N$_2$: \qty{2.9}{K},
जो केवल द्रव हीलियम के पास जमता है, अतः नाइट्रोजन सदा
$\tfrac52 R$ दिखाता है। (ग) चबूतरे $\tfrac32 R$ ($\sim\qty{50}{K}$ से नीचे),
$\tfrac52 R$ ($\sim200$ से $\sim\qty{1000}{K}$ तक), और
$\tfrac72 R$ की ओर चढ़ान, $\theta_{\text{कं}} \approx
\qty{6000}{K}$ के पास। (घ) विषम और सम $J$
भिन्न नाभिकीय चक्रण-जातियों के हैं, जो धीरे-धीरे आपस में बदलती हैं,
अतः ठंडे हाइड्रोजन का $C_V$ उसके ऑर्थो--पैरा इतिहास पर निर्भर करता
है --- नाभिकीय सांख्यिकी, किसी कैलोरीमापी से जाँची हुई।
\end{solution}

\begin{solution}{pb:b3:canonical-ensemble:1}
\textbf{1.} $V = \tfrac43\pi r^3 = \qty{4.0e-20}{m^3}$: $m' = V
\Delta\rho = \qty{8.3e-18}{kg}$, भार $m'g = \qty{8.1e-17}{N}$।
\textbf{2.} $n(h) = n_0\,\eu^{-m'gh/k_{\text{B}}T}$: अर्थात्
वायुदाबमापी नियम, जो सिकुड़कर सूक्ष्मदर्शी की स्लाइड पर आ गया।
\textbf{3.} $h_0 = k_{\text{B}}T/m'g = \num{4.04e-21}/\num{8.1e-17}
\approx \qty{50}{\micro m}$।
\textbf{4.} $h_0 \propto 1/r^3$: त्रिज्या में $10\%$ का फैलाव
पैमाना-ऊँचाई में $30\%$ का फैलाव है --- बहु-आकारी कणिकाएँ सीढ़ी को
लुगदी बना देती हैं। पेरँ ने महीनों बार-बार अपकेंद्रण करके छँटाई की।
\textbf{5.} वही सूत्र, पर द्रव्यमान $10^{10}$ गुना अलग: कणिकाओं का
``वायुमंडल'' $10^{10}$ गुना उथला है --- किलोमीटर सिकुड़कर दसियों
माइक्रोमीटर हो जाते हैं, और ठीक यही उसे सूक्ष्मदर्शी के नीचे पूरा का
पूरा देखने योग्य बनाता है।
\textbf{6.} लगातार अनुपात $0.55$, $0.55$, $0.57$: गिनती की त्रुटि
के भीतर अचर --- अर्थात् चरघातांकी। $h_0 =
\qty{30}{\micro m}/\ln(100/55) \approx \qty{50}{\micro m}$।
\textbf{7.} $k_{\text{B}} = m'g\,h_0/T = \num{8.1e-17} \times
\num{5.0e-5}/293 \approx \qty{1.4e-23}{J/K}$।
\textbf{8.} $N_{\text{A}} = R/k_{\text{B}} \approx
\num{6.0e23}\,\unit{mol^{-1}}$।
\textbf{9.} यहाँ (आदर्शीकृत आँकड़ों के साथ) कुछ ही प्रतिशत के भीतर;
पेरँ के असली प्रयोग आधुनिक मान के चारों ओर $\pm15\%$ बिखरे थे
--- हाथ से गिनी कणिकाओं के लिए चकित करने वाला, और कोटि के लिए
पूर्णतः निर्णायक।
\textbf{10.} $m_{\text{H}} = \qty{e-3}{kg/mol}/N_{\text{A}} =
\qty{1.7e-27}{kg}$।
\textbf{11.} स्थूल रसायन $R$ देता है; प्रकाशिक सूक्ष्मदर्शिकी और
तुलाई $m'$ देती हैं; और गिनती सीढ़ी: मेज़ पर के तीन मापन मिलकर उस
परमाणु का द्रव्यमान त्रिकोणित कर देते हैं, जिसे कोई देख नहीं सकता।
\textbf{12.} कुछ नहीं: व्युत्पत्ति ने केवल ``कोई तंत्र, जो किसी
भंडार से ऊर्जा का विनिमय करे'' का उपयोग किया --- बोल्ट्ज़मान का गुणक
आकार-अंधा है, और ठीक यही पेरँ ने सत्यापित किया।
\textbf{13.} $D = k_{\text{B}}T/6\pi\eta r =
\num{4.04e-21}/\num{4.0e-9} \approx \qty{e-12}{m^2/s}$।
\textbf{14.} प्रति मिनट $\sqrt{2Dt} = \sqrt{\num{1.2e-10}} \approx
\qty{11}{\micro m}$: नेत्रिका के अंशांकित जाल और धैर्य से आराम से
मापने योग्य।
\textbf{15.} दो स्वतंत्र परिघटनाएँ, दो स्वतंत्र सूत्र,
एक ही संख्या: स्थैतिक सीढ़ी और गतिक डगमगाहट की सहमति ने संयोग के लिए
कोई कोना नहीं छोड़ा --- आणविक परिकल्पना ने दोनों की भविष्यवाणी की थी।
\textbf{16.} प्रति घटक $\sqrt{k_{\text{B}}T/m} = \sqrt{\num{4.04e-21}/
\num{4.8e-17}} \approx \qty{9}{mm/s}$।
\textbf{17.} कणिका पर प्रति सेकंड $10^{19}$ चोटें पड़ती हैं और वह
नैनोमीटरों के भीतर अपना वेग भूल जाती है: प्रक्षेप्य उड़ान
अप्रेक्ष्य है, और आँख जो देखती है वह उसका समाकल है --- वही
विसरणी यादृच्छिक चाल।
\textbf{18.} स्टोक्स प्रतिकर्ष (वर्ष 2 के खंड के तरलों की श्यानता)
$6\pi\eta r$ देता है; और \omterm{thm:b3:canonical-ensemble:boltzmann}{विहित समुदाय}
$k_{\text{B}}T$: आइंस्टाइन का $D$ उन्हीं का भागफल है --- एक ही भिन्न
में यांत्रिकी और सांख्यिकी।
\textbf{19.} किसी कल्पना को गिना नहीं जा सकता: जैसे ही $N_{\text{A}}$
दो मापी गई संख्याओं का अनुपात बन गया, त्रुटि-सीमाओं समेत, परमाणु
प्रयोग की वस्तुएँ हो गए --- ऑस्टवाल्ड ने 1909 में छपकर हार मानी।
\textbf{20.} आकाश-नीला प्रकीर्णन, मापे गए इलेक्ट्रॉन-आवेश के साथ
विद्युत्-अपघटन, रेडियम से जमा हुआ हीलियम: चार असंबद्ध
भौतिक वाहिकाओं का एक ही $6\times10^{23}$ पर मिलना उस संख्या को किसी
सिद्धांत का नहीं, प्रकृति का गुण बना गया।
\textbf{21.} गुरुत्वीय ऊर्जा कणिकाओं को नीचे खींचती, विन्यास-एन्ट्रॉपी
उन्हें ऊपर फैलाती, और सौदा दर $T$ पर:
वह सीढ़ी $F = E - TS$ का न्यूनतम \emph{ही} है।
\textbf{22.} $h_0 \propto 1/g$: $10^5g$ पर प्रोटीन के लिए
$h_0 = k_{\text{B}}T/(m'\times10^5g) = \num{4.04e-21}/
\num{9.8e-17} \approx \qty{40}{\micro m}$ ---
अर्थात् अवसादन संतुलन, यानी पेरँ का प्रयोग, जो जैवरसायन विभागों में
रोज़ चलता है।
\textbf{23.} ``बहुत ऊँची'' ($h_0 \gg$ क्षेत्र की गहराई: कोई दृश्य
प्रवणता नहीं) से लेकर ``बहुत पतली'' ($h_0 \lesssim r$) तक: उपयोगी
त्रिज्याएँ मोटे तौर पर $0.1$--$\qty{0.5}{\micro m}$ फैली हैं ---
पेरँ का चुनाव भाग्य नहीं, अभिकल्पना था।
\textbf{24.} अब जब $k_{\text{B}}$ परिभाषा से ही यथार्थ है, वही
प्रयोग कणिकाओं का (उनके आकार या घनत्व का) अंशांकन करता है या
व्यवस्था की जाँच: भौतिकी --- बोल्ट्ज़मान की सीढ़ी --- अछूती है;
बस यह बदला है कि कौन-सी राशि अज्ञात मानी जाती है।
\textbf{25.} \qty{0.2}{\micro m} की किसी कणिका की जनसंख्या हर
$\sim\qty{35}{\micro m}$ पर आधी हो जाती है; और $n_0\eu^{-m'gh/k_{\text{B}}
T}$ से पढ़ने पर उस सीढ़ी ने $N_{\text{A}} \approx \num{6e23}$ लौटाया:
अवोगाद्रो की संख्या कणिका-दर-कणिका गिनी गई --- और परमाणु की बहस
सूक्ष्मदर्शी के मंच पर बंद हो गई।
\end{solution}
